\documentclass[../main.tex]{subfiles}
\ifSubfilesClassLoaded{\setcounter{chapter}{6}}{}
\begin{document}

\chapter{Organisasi Multi-berkas: \texttt{.c}, \texttt{.h}, dan Include Guards}

\begin{subcpmk}
  \item Sub-CPMK 3.2: Pemisahan \texttt{.c}/\texttt{.h}, \textit{include guards}, kompilasi multi-unit, \texttt{extern}
\end{subcpmk}

\SesiRPS{7}{3.2}{$\sim$170 menit}{CPMK-3}

% ============================================================
\section{Sarana dan prasarana}
% ============================================================
\texttt{gcc} dengan beberapa berkas sumber; editor yang mendukung banyak tab/panel.

% ============================================================
\section{Mengapa Multi-berkas?}
% ============================================================

Seiring pertumbuhan program, menyimpan semua kode dalam satu file \texttt{.c} menjadi tidak praktis. Program ribuan baris dalam satu file sulit dinavigasi, sulit dikelola oleh tim, dan memperlambat kompilasi karena setiap perubahan kecil mengharuskan seluruh file diproses ulang. Organisasi multi-berkas memisahkan \textbf{antarmuka} (deklarasi di \texttt{.h}) dari \textbf{implementasi} (definisi di \texttt{.c}), mirip daftar isi buku versus isi babnya.

Konsep ini merupakan fondasi pemrograman modular dalam C. Setiap ``modul'' terdiri dari sepasang file: header (\texttt{.h}) yang berisi prototipe fungsi dan deklarasi tipe, serta implementasi (\texttt{.c}) yang berisi kode fungsi. File \texttt{main.c} menjadi titik masuk program yang menggunakan modul-modul tersebut melalui \texttt{\#include}.

% ============================================================
\section{Deklarasi vs Definisi}
% ============================================================

{\small
\begin{tabularx}{\textwidth}{@{}|>{\raggedright\arraybackslash}p{2.2cm}|>{\raggedright\arraybackslash}p{5.8cm}|>{\raggedright\arraybackslash}X|@{}}
\hline
\textbf{Aspek} & \textbf{Deklarasi} & \textbf{Definisi} \\
\hline
Apa itu? & Memberitahu kompilator bahwa sesuatu \textit{ada} & Membuat sesuatu \textit{benar-benar ada} (alokasi memori / kode) \\
\hline
Di mana? & Header (\texttt{.h}) & Implementasi (\texttt{.c}) \\
\hline
Contoh fungsi & \texttt{int tambah(int, int);} & \texttt{int tambah(int a, int b)} \newline \texttt{\{ return a+b; \}} \\
\hline
Contoh variabel & \texttt{extern int total;} & \texttt{int total = 0;} \\
\hline
Boleh diulang? & Ya (di tiap file yang include) & Tidak (hanya satu file) \\
\end{tabularx}
\par
}

Setelah memahami perbedaan deklarasi dan definisi, berikut studi kasus praktis yang membangun proyek multi-berkas sederhana.

% ============================================================
\section{Studi Kasus: Modul \texttt{mathutil}}
% ============================================================

\subsection{File \texttt{mathutil.h} --- Header (antarmuka)}

\begin{lstlisting}[language=C,caption=\texttt{mathutil.h} dengan include guard]
#ifndef MATHUTIL_H    // jika belum didefinisikan
#define MATHUTIL_H    // tandai sudah didefinisikan

// Prototipe fungsi (deklarasi)
int tambah(int a, int b);
int kurang(int a, int b);
int kali(int a, int b);
double bagi(int a, int b);

// Konstanta
#define PI 3.14159265

#endif  // MATHUTIL_H
\end{lstlisting}

\subsection{File \texttt{mathutil.c} --- Implementasi}

\begin{lstlisting}[language=C,caption=\texttt{mathutil.c} --- definisi fungsi]
#include "mathutil.h"  // include header sendiri
#include <stdio.h>

int tambah(int a, int b) { return a + b; }
int kurang(int a, int b) { return a - b; }
int kali(int a, int b)   { return a * b; }

double bagi(int a, int b) {
    if (b == 0) {
        printf("Error: pembagian oleh nol!\n");
        return 0.0;
    }
    return (double)a / b;
}
\end{lstlisting}

\subsection{File \texttt{main.c} --- Program utama}

\begin{lstlisting}[language=C,caption=\texttt{main.c} menggunakan modul mathutil]
#include <stdio.h>
#include "mathutil.h"   // include header modul sendiri

int main(void) {
    int x = 10, y = 3;
    
    printf("%d + %d = %d\n", x, y, tambah(x, y));
    printf("%d - %d = %d\n", x, y, kurang(x, y));
    printf("%d * %d = %d\n", x, y, kali(x, y));
    printf("%d / %d = %.2f\n", x, y, bagi(x, y));
    printf("PI = %.5f\n", PI);
    
    return 0;
}
\end{lstlisting}

% ============================================================
\section{Include Guards}
% ============================================================

Include guard mencegah isi header disertakan lebih dari sekali dalam unit kompilasi yang sama. Tanpa include guard, jika dua header saling meng-\texttt{include} atau \texttt{main.c} secara tidak langsung menyertakan header yang sama dua kali, kompilator akan menemukan deklarasi duplikat dan menghasilkan error.

Ada dua cara menulis include guard:

\begin{modultable}[]{Perbandingan include guard tradisional dan \texttt{\#pragma once}}
\begin{tabular}{|l|p{5cm}|p{5cm}|}
\hline
\textbf{Metode} & \textbf{Tradisional (\texttt{\#ifndef})} & \textbf{Modern (\texttt{\#pragma once})} \\
\hline
Sintaks &
\texttt{\#ifndef HEADER\_H} \newline
\texttt{\#define HEADER\_H} \newline
\texttt{...isi...} \newline
\texttt{\#endif} &
\texttt{\#pragma once} \newline
\texttt{...isi...} \\
\hline
Portabilitas & Standar C (semua kompilator) & Didukung mayoritas kompilator (bukan standar resmi) \\
\hline
\end{tabular}
\end{modultable}

% ============================================================
\section{Kompilasi Multi-berkas}
% ============================================================

\begin{lstlisting}[language=bash,caption=Langkah kompilasi terpisah]
# Langkah 1: kompilasi tiap .c menjadi .o (object file)
gcc -Wall -Wextra -std=c17 -c mathutil.c -o mathutil.o
gcc -Wall -Wextra -std=c17 -c main.c -o main.o

# Langkah 2: tautkan (link) semua .o menjadi executable
gcc mathutil.o main.o -o program

# Atau satu baris (gcc langsung):
gcc -Wall -Wextra -std=c17 -o program main.c mathutil.c
\end{lstlisting}

\begin{konsep}
\textbf{Keuntungan kompilasi terpisah:} jika hanya \texttt{mathutil.c} yang diubah, cukup kompilasi ulang file itu saja, lalu tautkan dengan \texttt{main.o} yang sudah ada. Pada proyek besar, ini menghemat waktu kompilasi secara signifikan.
\end{konsep}

% ============================================================
\section{Keyword \texttt{extern}}
% ============================================================

Kata kunci \texttt{extern} digunakan untuk mendeklarasikan bahwa sebuah variabel \textbf{didefinisikan di file lain}. Ini memungkinkan berbagi variabel global antar file --- namun praktik ini harus diminimalisir (lebih baik menggunakan parameter fungsi).

\begin{lstlisting}[language=C,caption=Penggunaan \texttt{extern}]
// config.c
int jumlah_mahasiswa = 40;  // DEFINISI (alokasi memori)

// main.c
extern int jumlah_mahasiswa;  // DEKLARASI (referensi ke config.c)
printf("Jumlah: %d\n", jumlah_mahasiswa);
\end{lstlisting}

% ============================================================
\section{Referensi sesi}
% ============================================================
\begin{itemize}[leftmargin=*]
  \item GNU C Reference Manual, praprosesor. \url{https://www.gnu.org/software/gnu-c-manual/gnu-c-manual.html}
  \item Beej's Guide, \textit{Multifile Projects}. \url{https://beej.us/guide/bgc/html/index.html}
  \item cppreference, \textit{include}. \url{https://en.cppreference.com/w/c/preprocessor/include}
\end{itemize}

% ============================================================
\begin{aktivitas}
  \item Bangun mini-proyek dua unit: utilitas (\texttt{.h} + \texttt{.c}) dan \texttt{main.c}.
  \item Kompilasi secara terpisah (\texttt{-c} untuk setiap \texttt{.c}, lalu tautkan).
  \item Coba hilangkan include guard dan \texttt{\#include} header dua kali --- amati error-nya.
\end{aktivitas}

\begin{latihan}
  \item Apa perbedaan \texttt{\#include <stdio.h>} dan \texttt{\#include "mathutil.h"}?
  \item Mengapa definisi fungsi tidak boleh diletakkan di file header?
  \item Jelaskan kegunaan \texttt{extern} dan risikonya.
\end{latihan}

\begin{asesmen}
\textbf{Tugas M5b:} struktur folder (minimal 3 file: header, implementasi, main), header aman dengan include guard, skrip kompilasi manual (perintah dalam \texttt{README}).
\end{asesmen}

\begin{checklist}
  \item Saya memahami perbedaan deklarasi di \texttt{.h} dan definisi di \texttt{.c}
  \item Saya dapat menulis include guard yang benar
  \item Saya dapat mengompilasi dan menautkan beberapa file \texttt{.c}
  \item Saya memahami kapan menggunakan \texttt{extern}
\end{checklist}

\begin{rangkuman}
Bab ini membahas organisasi program C menjadi beberapa file: header (\texttt{.h}) untuk deklarasi dan implementasi (\texttt{.c}) untuk definisi. Include guard mencegah penyertaan ganda. Kompilasi terpisah meningkatkan efisiensi pada proyek besar.

\textbf{Poin kunci:} include guards (\texttt{\#ifndef}); model kompilasi terpisah lalu taut; \texttt{"..."} untuk header lokal, \texttt{<...>} untuk header sistem.

\textbf{Kata kunci:} header, \texttt{.h}, \texttt{.c}, include guard, \texttt{extern}, kompilasi terpisah, linking
\end{rangkuman}

\end{document}
